\section{Conclusion}

\subsection{Achievement of Objectives}

All objectives stated in section 7.2.1 of the laboratory manual have been met. The application implements a complete, table-driven Moore finite state machine with two states (\texttt{LED\_OFF} and \texttt{LED\_ON}) that toggles on every confirmed press of a push button. The transitions are governed by a constant table that mirrors one-to-one Tabelul~7.1 of the manual, so the formal model and the source code are in lock-step. A dedicated \texttt{Button} driver performs software debouncing with a 30~ms stable-time filter and exposes a clean rising-edge API to the FSM, which keeps the application loop short and free of timing-related branches. The current FSM state is published in real time on the serial terminal through the STDIO library (\texttt{printf} redirected onto UART0), satisfying the manual's reporting requirement.

The supplementary requirements of the chapter ("Cerințe suplimentare") are also honoured. The FSM logic lives in a dedicated software module, separable from the input and output drivers. The behaviour is documented by both a state diagram and an algorithmic flowchart. The debounce is implemented in software, with no extra hardware. All low-level peripheral access is encapsulated inside the four reusable libraries under \texttt{labs/lib/}: the lab entry point itself contains no \texttt{digitalRead}, \texttt{digitalWrite}, \texttt{pinMode}, or \texttt{Serial.*} call --- a strict reading of the project rule that no low-level code may live outside libraries.

\subsection{Performance Analysis and System Limitations}

The system performs comfortably within the constraints set in the requirements. The worst-case latency between a press and the corresponding state report is approximately 35~ms (30~ms debounce window plus 5~ms loop tick), which is well below the 100~ms budget. The firmware footprint is small (about 5.9~KB of Flash and 833~bytes of SRAM), leaving more than 95~\% of the on-chip resources free for future extensions.

The current implementation has two notable limitations. First, the main loop uses a blocking \texttt{delay(5)} between iterations, which is acceptable for this single-task application but would prevent the same code from coexisting with another time-sensitive activity. Second, the Moore output is bound to a single LED; supporting multiple actuators or LCD-based reporting would require minor wiring changes at the application layer. Neither limitation affects the correctness of the FSM, only its scope.

\subsection{Proposed Improvements}

Several extensions are natural follow-ups. The blocking delay could be replaced by the cooperative \texttt{TaskScheduler} module already used in Lab 2.1, allowing the application to run additional periodic tasks (e.g., LED heartbeat, statistics) alongside the FSM. The reporting could be moved from the serial terminal to an LCD via the \texttt{LcdDisplay} module already available in \texttt{labs/lib/}, which would only require a new application-layer adapter and no change in the FSM or in the \texttt{Button} driver. The FSM itself could be extended with intermediate states (e.g., a \texttt{LED\_BLINK} state activated by a long press) to demonstrate state-dependent timed behaviour and to introduce the Mealy variant. Finally, replacing the polled \texttt{Button} with an interrupt-driven \texttt{INT}-pin handler would shrink the worst-case latency below the debounce window itself, with the trade-off of slightly more complex synchronisation.

\subsection{Reflections on the Learning Experience}

Translating the textbook description of a Moore automaton into a constant table that lives at the heart of the firmware proved to be the most valuable insight of this lab. The 4-line state table in \texttt{ButtonLedFsm.cpp} is not just an implementation detail: it is the formal specification of the controller and is the same artefact that appears in the manual, in the design diagram, and in the code review. When the diagram, the table, and the implementation must agree by construction, ambiguity is impossible and changes are localised. Equally instructive was the discipline of keeping all hardware access out of the lab entry point: the resulting application code reads almost like pseudocode, which makes inspection, code review, and testing markedly easier than in earlier labs where small \texttt{digitalWrite} fragments were sprinkled across modules.

\subsection{Impact of Technology in Real-World Applications}

The pattern realised in Lab 6.1 --- a debounced discrete input feeding a Moore FSM whose output drives a stable physical actuator --- is foundational in industrial and consumer embedded systems. Power buttons of laptops and smartphones, smart home wall switches that toggle relays, automotive cabin functions such as cruise-control engagement or hazard lights, industrial safety interlocks, and avionics control panels are all designed around the same idea. In safety-critical contexts, the explicit state table is invaluable for certification: certifying authorities (DO-178C in avionics, ISO~26262 in automotive, IEC~62304 in medical) require traceability between requirements, design, and implementation, and a table-driven FSM provides that traceability almost for free. The same modular discipline applied here --- isolating FSM logic, separating driver code, and using a portable C standard I/O abstraction --- is the standard practice codified in industry frameworks such as AUTOSAR.
